\section{The squeeze between two slopes}\label{sec:lattices}

This section describes how the inequality $\ASK\le\MAIN+\Err$ is proved. The objects are
Hermitian lattices over the ring of integers of a number field, and the tool is the slope
method of Bost \cite{bost1996,bost2001}, which rests on the filtration of Harder and
Narasimhan \cite{harder1975} in the arithmetic form due to Stuhler and Grayson
\cite{stuhler1976,grayson1984}.

\subsection{Slopes}
Let $K$ be a totally complex number field. A \emph{Hermitian lattice} over $\cO_K$ is a finitely
generated projective $\cO_K$-module $L$ with a positive Hermitian form on $L\otimes_v\C$ for
each infinite place $v$. Its slope $\mu(L)$ is its Arakelov degree divided by its rank, in the
normalisation \lvS
\[
\mu(L)=2\,\frac{\widehat{\deg}\,L}{[K:\Q]\,\rk L}-\log 2+\frac{\log\abs{d_K}}{[K:\Q]},
\]
and $\mu_{\max}$, $\mu_{\min}$ denote the largest and smallest slopes of the pieces of its
Harder--Narasimhan filtration. The height $\htt(\varphi)$ of a $K$-linear map between such
lattices is the mean over the infinite places of $\log\norm{\varphi}_v^2$.

\subsection{The cast}\label{sec:cast}
Fix a translate $u_0$. The index set is a finite set of ``words'' $u$; each word determines
an archimedean type of automorphic form. For each $u$ there are \lvS:
\begin{itemize}
\item a lattice $M_u$: an integral structure on a space of vector-valued algebraic modular
  forms, in the sense of Gross \cite{gross1999}, on a definite unitary group over the host
  field $E$, with the Petersson norm at one place. The space is the whole space of forms of
  the given weight and level, namely all functions on the finite set of classes of a genus
  with values in the weight representation, invariant under the stabilisers
  (\lean{Fi}); nothing is selected by Hecke eigenvalues;
\item a sublattice of $M_u$ of forms ``ramified at $T$'', where $T=T(u)$ is an auxiliary level
  prime, and its rescaling
  \begin{equation}\label{eq:Nu}
  N_u\;=\;\Bigl(\prod_{\lambda}\fP_\lambda^{-\lfloor e_\lambda\,g(u,\lambda)\rfloor}\Bigr)\cdot(\text{that sublattice}),
  \end{equation}
  where $\lambda$ runs over finite places of the coefficient field, $e_\lambda$ is a ramification
  index, and the exponent is a sum over the $m$ lifts,
  \[
  g(u,\lambda)=\sum_{j=1}^{m}\bigl(\mathrm{gain}_j(\lambda)+\mathrm{fo}_j(\lambda)-\mathrm{ee}_j(u,\lambda)\bigr),
  \]
  of the \emph{gain} at $\lambda$ for the $j$-th lift (\cref{sec:adm}) and of two correction
  terms, of which the second depends on the word (\lean{Iu}, \lean{NprimeDef}, \lean{sg}).
  The record does not make $g(u,\lambda)$ non-negative;
\item a map $B_u$ from the space of $M_u$ to a fixed space $W$, sending a form to the periods
  of its $m$-fold iterated theta lift against test vectors;
\item a real number $h(u)$ bounding the height of $B_u$.
\end{itemize}
There is one target: a lattice $P\subset W$, the dual of a lattice $R$ glued prime by prime
from images of integral forms under the lift. We call $P$ the \emph{lattice of periods}.

The final theorem requires, for the family at each scale, a record of eleven clauses
(\lean{SizeData}\textsubscript{\ttfamily 2}) \lvS. Four of them make the chain:
\begin{align}
&\ASK(u_0)-\mathrm{Ga}(u_0)\ \le\ \mathrm{Lo}, \tag{\textsf{Loss}}\label{eq:loss}\\
&\mu(M_u)+\mathrm{Ga}(u_0)+\mathrm{Fo}-\mathrm{Ee}(u)-\mathrm{Lo}_5(u)\ \le\ \mu(N_u),
   \tag{\textsf{Gain}}\label{eq:gain}\\
&\bar\mu-\log\abs I-\frac{1}{1-\delta_F}\Bigl(4u_F\log\bigl(1+\tfrac{HH}{u_F}\bigr)
     +(\eta_F+4\sqrt{e_F})\,HH\Bigr)\ \le\ \mu(P), \tag{\textsf{Slope}}\label{eq:slope}\\
&\mu(P)+h(u)-\mathrm{Fo}\ \le\ \sum_{t\in TT}\Bigl(m\bigl(\ph(a_t)-a_t\log\Omega\bigr)-a_t\log T(u)\Bigr)
     +(\text{small}). \tag{\textsf{Div}}\label{eq:div}
\end{align}
Here a sum over $\lambda$ weighted by $2e_\lambda$ times a normalised logarithm of the norm of
$\lambda$ turns the three parts of the exponent into three real numbers: $\mathrm{Ga}(u_0)$,
the total gain, from which a cost $\tau_\lambda$ at each place is first subtracted;
$\mathrm{Fo}$; and $\mathrm{Ee}(u)$ (\lean{Ga}, \lean{Fo}, \lean{Ee}). Further, $\bar\mu$ is the
rank-weighted mean over $u$ of $\mu(N_u)-h(u)$; the term marked small is an error of the kind
of \cref{sec:balance}, with a factor $x^{e_D\eps}$; and $\mathrm{Lo}$, $\mathrm{Lo}_5(u)$ are
losses that are small. Read in order: the count is at most the total gain \eqref{eq:loss};
rescaling raises the slope of $N_u$ above that of $M_u$ by the total gain and the two
corrections \eqref{eq:gain}; the mean slope of the $N_u$, less heights, is at most the slope
of $P$ \eqref{eq:slope}; and the slope of $P$ is at most the period term on the right of
\eqref{eq:div}. The term $\mathrm{Fo}$ enters \eqref{eq:gain} and \eqref{eq:div} with opposite
signs and cancels. The term $\mathrm{Ee}(u)$ is not small, and it is cancelled otherwise, as
follows.

Four further clauses close the chain. One bounds $\mu(M_u)$ from below and
$\mu_{\max}(M_u)$ from above by error terms (\lean{Slopes}); one bounds the variation with $u$
of $h(u)+\sum_ta_t\log T(u)$ by error terms (\lean{HComp}); one bounds $\log T(u)$ from
below by a dimension count (\lean{Count5}, \cref{sec:period}); and one (\lean{EComp5}) says
that $\mathrm{Ee}(u)$ equals $m\,\ell_g\sum_ta_t$ up to an error term, where
$\ell_g$ is the logarithm of the product of two auxiliary level primes of the host. The lower bound
for $\log T(u)$ contains the term $m\,\ell_g$, so after the dimension count the right side of
\eqref{eq:div} is at most $\MAIN-m\,\ell_g\sum_ta_t$ up to an error term, by an exact identity
(\cref{st:key}); and the term $m\,\ell_g\sum_ta_t$, which is not an error term,
cancels $-\mathrm{Ee}(u)$ in \eqref{eq:gain}. The remaining three clauses are positivity and
size conditions.

\subsection{Why $N_u$ still maps into $P$}\label{sec:adm}
Rescaling the sublattice by negative powers of primes raises its slope by the total gain and
the two corrections, less what passing to the sublattice costs, but it is only useful if the
rescaled lattice is still
mapped into $P$ by $B_u$. This is the statement
\lvS{} (\lean{adm2}):
\begin{quote}
for every form $\varphi$ in the unscaled sublattice, every place $\lambda$ and every integral test
vector $F$, the period of the lift of $\varphi$ against $F$ is $0$ or has $\lambda$-adic valuation at
least $\lfloor e_\lambda\,g(u,\lambda)\rfloor$.
\end{quote}
It is proved by \lean{admissible_level_two}, and it is where the arithmetic of
\cref{sec:divisibility} is spent.

The gain is a closed formula \lvS{} (\lean{gainP}), one for each of the $m$ lifts. It is used
at the primes of two of the classes into which the host sorts the primes, the inert class and
one other; at the rest the gain is $0$, and it is replaced by $0$ at a set of dropped primes,
which contains every prime of the inert class for which $h_p$ below has odd order and is not
$0$. For a prime $p$ let $h_p\in H$ be the
component in $H$ of its Frobenius element. The top places are the elements $t$ of $TT$; they
fall into orbits under translation by $h_p$ (by $2h_p$ for the primes of one class), and each
orbit has a fixed base point. Call a
place of \emph{even position} if the least number of steps of $h_p$ that leads to it from the
base point of its orbit is even (\lean{IsV}; positions are used only at primes of the inert
class). For a top place $t$ and $e\ge1$ the
\emph{partner} is the place
$t+e\,h_p$ (taken on a second copy of the places when $p$ is of that class and $e$ is odd), and
the \emph{capacity} $c_t(e)$ is the weight at the partner less an amount that depends on the
lift (\lean{capU}). On one orbit, with sizes $\alpha_t$ equal to $a_t$ less a fixed constant, put
\begin{align}
g_U&=\gamma_p\sum_{t}\alpha_t\;-\;\sum_{t\ \text{of odd position}}\ \sum_{e\ \mathrm{odd}}
\min\Bigl(\frac{\alpha_t}{p^{e}},\,c_t(e)\Bigr)\;-\;\delta_p,
\qquad \gamma_p=\frac{p}{p^2-1}=\sum_{e\ \mathrm{odd}}\frac1{p^{e}},\label{eq:gU}\\
g_0&=\sum_{t}\ \sum_{e\ \mathrm{odd},\ \text{not excluded}}\max\Bigl(\frac{\alpha_t}{p^{e}}-c_t(e),\,0\Bigr)-\delta_p ,\notag
\end{align}
the sums over the top places $t$ of the orbit (\lean{gU}, \lean{g0}). Here $\delta_p$ is a
correction of lower order that depends on the orbit, and an odd $e$ is \emph{excluded} for $t$
if the partner of $t$ at $e$ is another top place, on the first copy, and $e$ is the least
positive exponent that reaches it (\lean{Excl}). The formula gives $0$ if $p$ is of the inert
class with $h_p=0$, or if $p$ exceeds a bound. Otherwise it gives
\[
\sum_{\text{orbits in }VV}\max\bigl(g_U-\mathrm{Pr},\,g_0,\,0\bigr)
\;+\;\sum_{\text{other orbits}}\max\bigl(g_0,\,0\bigr),
\]
with a penalty $\mathrm{Pr}$. The set $VV$ of orbits (\lean{VV}) is empty unless $p$ is a
\emph{candidate}, that is a prime of the inert class below a bound whose orbits have length at
least $x^{\eps/4}$ and which, if it is above a second bound, has no power within a small
ratio of the first, fixed at each scale, and a smallness condition holds; then it is the
set of orbits that satisfy a regularity condition on the positions of their top places
(\lean{Good}).

Since $\gamma_p\alpha_t=\sum_{e\,\mathrm{odd}}\alpha_t/p^e$, a place of odd position
contributes to $g_U$, as it does to $g_0$ at the exponents that are not excluded,
\[
\sum_{e\,\mathrm{odd}}\max\bigl(\alpha_t/p^e-c_t(e),\,0\bigr),
\]
and a place of even position contributes the whole of $\gamma_p\alpha_t$. This
is the shape of the count \eqref{eq:ask}, in which only odd powers of Frobenius occur because
$\Phi$ has index two, with the capacity in the place of the weight. That is why
\eqref{eq:loss} holds with a small loss. The divisibility of \cref{sec:divisibility} supplies
the rate $\gamma_p$.

\subsection{The slope inequality}\label{sec:slope}
The clause \eqref{eq:slope} is an instance of a theorem about Hermitian lattices that involves
no automorphic forms. We state it in full because it is the cleanest self-contained new
statement we met.

\begin{statement}[\lean{TheoremF}]\label{st:F}
Let $K$ be a totally complex number field, $P$ a Hermitian lattice over $\cO_K$ in a $K$-space
$W$, and for $i$ in a finite set $I$ let $N_i$ be a Hermitian lattice in a $K$-space $V_i$ and
$B_i\colon V_i\to W$ a $K$-linear map with $B_i(N_i)\subseteq P$ and $\htt(B_i)\le h_i$ whenever
$B_i\ne0$. Let $\bar\mu\in\R$, $u_F\ge0$ and $HH\ge0$ satisfy
\[
\bar\mu\sum_i\rk N_i\le\sum_i\rk N_i\,\bigl(\mu(N_i)-h_i\bigr),\qquad
\mu_{\max}(N_i)-h_i\le\bar\mu+u_F,\qquad \bar\mu-\mu_{\min}(P)\le HH .
\]
Suppose that at one infinite place $v$, after a change of bases, after extending $W_v$ to a
space $W_v'$ in which the added part has dimension at most a fraction $\delta_F<1$ of the
total, and after extending each $B_i$ to a map into $W_v'$, arbitrarily on the added part,
there are weights $\gamma_i>0$
with
\begin{gather*}
\sum_i\norm{\gamma_iB_i^*B_i-1}_{\mathrm{HS}}^2\le e_F\sum_i\rk N_i,\qquad e_F\le\tfrac14,\\
\tr(Q^2)\le(1+\eta_F)\,\frac{(\tr Q)^2}{\dim W_v'},\qquad Q=\sum_i\gamma_iB_iB_i^* .
\end{gather*}
Then
\[
\mu(P)\ \ge\ \bar\mu-\log\abs I-\frac1{1-\delta_F}\Bigl(4u_F\log\bigl(1+\tfrac{HH}{u_F}\bigr)+(\eta_F+4\sqrt{e_F})\,HH\Bigr),
\]
the term $4u_F\log(1+HH/u_F)$ being read as $0$ when $u_F=0$.
\end{statement}

The classical slope inequality says that an injective map $\varphi\colon N\to P$ of height $h$
forces $\mu_{\max}(N)\le\mu_{\max}(P)+h$; it bounds the top of the target. \Cref{st:F} bounds
the slope of the whole target from below, from many maps at once. The hypothesis on $Q$ says
that the maps form a nearly tight frame: each is nearly an isometry up to a scalar, and their
images together nearly fill the target evenly.

The proof runs level
by level. For each $\tau\in[0,HH]$: cut each $N_i$ at the step of its Harder--Narasimhan
filtration of slope $\bar\mu-\tau+h_i$; a Markov inequality (\lean{frameB_share}) shows that the parts
cut away have at most the share $u_F/(u_F+\tau)$ of the total rank; the classical inequality, in
its $\mu_{\min}$ form, sends what is kept into the step of the filtration of $P$ of slope
$\bar\mu-\tau$; and a lemma of linear algebra on nearly tight frames shows that the images of what is
kept span all but a share $\eta_F+4u_F/(u_F+\tau)+4\sqrt{e_F}$ of the frame space. So the filtration of
$P$ at slope $\bar\mu-\tau$ has almost full rank for every $\tau$, and integrating over $\tau$ gives
the bound: $\int4u_F/(u_F+\tau)\,d\tau$ is the term $4u_F\log(1+HH/u_F)$ and $\int(\eta_F+4\sqrt{e_F})\,d\tau$
is the term $(\eta_F+4\sqrt{e_F})\,HH$. The term $\log\abs I$ is a slack: the form of the theorem
without it is also proved, and the form with it is obtained by discarding a non-negative
quantity.

\subsection{The period side}\label{sec:period}
The clause \eqref{eq:div} bounds $\mu(P)$ from above. Three things enter.

\emph{A determinant.} The $m$-fold lift is a composite of $m$ single lifts
$\theta_1,\dots,\theta_m$ through auxiliary spaces. A general inequality about Hermitian
lattices over a CM field \lvS{} (\lean{BoundsMu}) takes as its hypothesis a factorisation of
determinants at every archimedean place, of the shape
$\det(\text{composite})=\kappa^{D}\prod_j\det(\theta_j)$, and bounds the slope of $P$ in terms
of the slopes of $m$ auxiliary lattices $M'_1,\dots,M'_m$, one for each lift. The scalars are
bounded below by a theorem \lvS{} (\lean{lamLow_of_hypotheses_alt_alt}) whose hypotheses are
seven statements of the literature, each proved in the archive: rapid decay of cusp forms,
the Siegel--Weil formula for unitary groups in two forms
\cite{weil1965,ichino2004,ichino2007}, the basic identity of the doubling method
\cite{psrallis1987} with the continuation of the Siegel Eisenstein series (under a name that
refers to the local zeta integrals; see \cref{sec:literature}), the unramified computation of
the local doubling integral, a finiteness theorem of Bernstein and Zelevinsky, and
Harish-Chandra's formula for the formal degree of a discrete series representation of a
unitary group. The identities have
the shape of the Rallis inner product formula \cite{rallis1984}.

\emph{A ledger.} The clause \eqref{eq:div} is assembled from one inequality per lift $j$ and
place $\rho$ \lvP{} (\lean{divisor_inequality_assembly}): a covolume term
$w_j(\rho)\log\Omega$, plus a kernel term $\ph(\nu)$, minus an eigenvalue term and minus
$\nu\log\Omega$, is at most $\ph(a_\rho)-a_\rho\log\Omega$ if $\rho$ is a top place and $0$
otherwise, up to errors. Here $\ph(\nu)\approx\log\nu!$ and $\Omega$ is a period. Summed over
the $m$ lifts and all places, and with $-a_t\log T$ from the last map, this is
$m\sum_t\bigl(\ph(a_t)-a_t\log\Omega\bigr)-\sum_ta_t\log T$, the right side of \eqref{eq:div}.
Every factorial but the top one has cancelled against an eigenvalue, and every
$\log\Omega$ but the top one against a covolume.

\emph{Arithmetic geometry.} The covolume of the auxiliary lattices is an input of a different
kind, and a large one. $M'_j$ is a lattice of integral forms on a moduli space attached to an
\emph{indefinite} unitary group, integrality meaning non-negative valuation at every point over
every finite place \lvS{} (\lean{cuspA}). Its covolume is bounded through a chain of slope
comparisons \lvS{} (\lean{lcv_le}): the covolume is at most minus the slope of a ``good''
sublattice cut out by an orthogonal idempotent; the complementary ``bad'' part has rank share
at most $3\exp(-x^{\eps_z/7})$, with $\eps_z$ an exponent of the scale; the slope of the lattice
of that type is at least a core slope
less a small error; and the core slope is at least minus a main term, which is within a small
error of $(p'-q')\sum_\tau\mathrm{anchor}_\tau\,\log\Omega$, with two integers $p'$, $q'$ of the
datum of the indefinite group and a sum over an index $\tau$ of that datum. The theorem that assembles this
\lvS{} (\lean{topCovolSlim2_use_alt_alt}) has fifteen hypotheses. Among them are resolution of
singularities and
principalisation in characteristic zero \cite{hironaka1964}, an arithmetic Hilbert--Samuel
statement \cite{gilletsoule1992,zhang1995,abbesbouche1995}, a bigness statement of Moriwaki,
toroidal compactifications \cite{faltingschai1990,lan2013}, Riemann--Roch for nef bundles,
semi-negativity of the Lie algebra of a semi-abelian scheme, Zarhin's trick, a statement of
Fargues, and Grothendieck's deformation theory; one hypothesis is a bundle of twenty-eight
further statements, among them heights of Chow forms and the vanishing theorems of
Kawamata--Viehweg and Nadel. So the slope argument
is not lattice theory and automorphic forms alone: it rests on arithmetic intersection theory
on a Shimura variety. Each of these statements is proved inside the archive.

\emph{The period.} $\Omega$ is a CM period. What the argument needs is its mean, and here the
archive proves the Chowla--Selberg formula
\cite{chowla1949epstein,selbergchowla1967} in the form \lvS{} (\lean{chowla_selberg_eta})
\begin{equation}\label{eq:cs}
\frac1{h_K}\sum_{c\in\mathrm{Cl}(K)}\log\bigl(4\pi\,\im\tau_c\,\abs{\eta(\tau_c)}^4\bigr)
=\gamma+\log\frac{2\pi}{\sqrt d}-\frac{L'}{L}(1,\chi_d),
\end{equation}
for an imaginary quadratic field $K$ of discriminant $-d$ with $d>4$ and its quadratic character
$\chi_d$, from the factorisation $\zeta_K=\zeta\cdot L(\cdot,\chi_d)$ and
Kronecker's first limit formula \cite{siegel1980advanced}, which is also proved there \lvS. This is how
$\tfrac{L'}{L}(1,\chi_\Phi)$ enters \eqref{eq:main}.

\emph{The count for $T$.} The dimension of the space of forms is known exactly: the number of
classes is a mass times an index times $(T^k-1)^f$ \lvS{} (\lean{ClassNumberFormula}\textsubscript{\ttfamily 2}). A
lower bound for $\log T(u)$ is the clause \lvS{}
(\lean{Count5}),
\[
\sum_{j=1}^m\frac1f\sum_{\rho}\log\frac{w_j(\rho)\,v}{\pi}+m\,\ell_g-c\Bigl(L_a+\frac{m^2L_a}{k}\Bigr)\ \le\ \log T(u),
\]
where $f=\abs\Phi$, the numbers $v$ and $\ell_g$ are constants of the data, $L_a$ is a
logarithmic size parameter and $c$ a constant of the record.

\begin{statement}[\lean{key_identity5}]\label{st:key}
With the notation above, the sum $m\sum_t\bigl(\ph(a_t)-a_t\log\Omega\bigr)$ on the right of
\eqref{eq:div}, less $\sum_ta_t$ times the lower bound for $\log T$ just displayed, plus
$m\,\ell_g\sum_ta_t$, equals
\[
\MAIN(u_0)\;+\;\Bigl(\sum_ta_t\Bigr)\,c\,\Bigl(L_a+\frac{m^2L_a}{k}\Bigr).
\]
\end{statement}
The $\log(\pi/v)$ inside $\log\Omega$ cancels the $\log(v/\pi)$ of the count. There is no
discrepancy term: the lattice side returns the analytic main term \eqref{eq:main} identically.
What is left is the term $m\,\ell_g\sum_ta_t$, which the clause \lean{EComp5} sets against
$\mathrm{Ee}(u)$ (\cref{sec:cast}), and an error of the shape of the first term of the bracket
of \cref{sec:balance}, with the same factor $1+m^2/k$.

\subsection{Chaining}
From \eqref{eq:loss}, \eqref{eq:gain}, \eqref{eq:slope}, \eqref{eq:div} and \cref{st:key},
with the four clauses that bound $\mu(M_u)$, the variation of $h(u)$, $\log T(u)$ and the
difference between $\mathrm{Ee}(u)$ and $m\,\ell_g\sum_ta_t$, for every translate $u_0$,
\[
\ASK(u_0)\ \le\ \MAIN(u_0)+\Err ,
\]
where $\ASK(u_0)$ is the sum over $t\in TT$ and $j\le m$ of the counts \eqref{eq:ask} with the
weights $w_{u_0,j}$, and $\MAIN(u_0)$ the corresponding sum of main terms. This is what
\cref{sec:reduction} needs. The archive proves it as
\lean{each_unit_bound_statement} \lvM, for every translate and not only on average.
